\chapter{Security and Robustness Analysis}
\section{Attack Surface Analysis}
\subsection{Payload Removal or Overwrite}
An actor with permission to alter the storage can remove a named stream,
replace its contents, or overwrite the bytes holding the slack marker. A
slack marker also depends on allocation remaining compatible with the
recorded location. The current implementation does not reserve the slack
through an active filesystem service and does not provide an error-correcting
code. Once the selected original bytes are lost, knowing their earlier hash
does not reconstruct them.

The single-bit experiments demonstrate a narrow consequence: recovery still
exports the requested length, but original-payload verification fails. They
do not establish a measured survival rate after mounted filesystem writes,
defragmentation, backup, copying to another filesystem, or secure deletion.
Those transformations would require new controlled images and independent
recording of their effects. For current evidence, stable offline storage is
an assumption, not a robustness feature.

\subsection{Manifest Substitution and Misleading Expectations}
The original hash is a trust anchor. If an attacker replaces both a payload
and the manifest used to judge it, matching those two values does not
establish agreement with the earlier legitimate content. SHA-256 digests
identify changes relative to a reference \cite{sha}; they are not signatures
proving the reference's author. The project's plaintext JSON manifests are
not cryptographically authenticated.

This attack is feasible wherever a person can edit both the evidence copy
and its expectations. Repository history offers a review trail, but an
investigator still needs to establish which version is trusted. A useful
hardening step would be to sign acquisition manifests and store the original
digest record separately. That would address provenance and substitution;
it would not hide the marker or improve NTFS parsing coverage.

\subsection{Malformed Metadata and Resource Exhaustion}
A hostile image can supply inconsistent lengths, invalid mapping pairs,
out-of-volume extents, unsupported flags, or damaged update sequences. These
structures can lead to false content, excessive work, or unexpected reads if
interpreted without checks. The recovery layer validates many such cases and
the suite includes malformed and truncated inputs. Unsupported layouts are
refused instead of treated as successful empty files.

The implementation has not undergone comprehensive fuzzing, formal
verification, or independent adversarial review. Some valid-looking metadata
can still request substantial work, and a full-image hash costs a read of
the entire input. Controlled fixtures do not establish resistance to every
resource-exhaustion attack. Operational hardening would include explicit
resource budgets, controlled process execution, and a documented policy for
very large images.

\subsection{Source Changes During Examination}
Before/after hashing is an observation over two points in time. It does not
hold a transactionally immutable snapshot throughout the export. If another
process changes the source while the tool is running, reads can observe an
inconsistent state. The code warns against changing the image between
location and recovery; the report can mark a changed-source hash as failure.
An adversary able to change and restore a file between checks is outside the
protection offered by those two hashes alone.

The appropriate mitigation is acquisition and stable-image handling, with
read-only media or a protected snapshot where available. A reproducible
student experiment should not be described as a hardware write-blocking
system or a complete chain-of-custody solution. NIST's forensic guidance
provides the wider process context for preserving and interpreting evidence
\cite{nist86}.

\section{Detection and Confidentiality Limits}
The markers are plain ASCII and the stream name is documented. The design
does not attempt to defeat an examiner who enumerates named streams or reads
the relevant slack bytes. The reviewed Streams and Sleuth Kit capabilities
illustrate why omission from a default file view is not equivalent to
undetectability \cite{sysstreams,icat}. Our results should therefore be read
as evidence that such material can be recovered, not that the hiding method
evades capable forensic tools.

No encryption protects the hidden content. Microsoft's EFS is a distinct
filesystem encryption feature \cite{efs}, and it is outside this extractor's
supported stream types. Adding application-level encryption in a future
experiment would address confidentiality of the payload, but would not make
its existence inherently invisible. Verification of plaintext recovery would
then require a clearly specified decryption stage and key-handling procedure.

\section{Implementation and Evaluation Limitations}
The release's recovery path applies fixups before interpreting resident
content. The older metadata CLI does not apply those checks on every path;
its output must not be treated as equivalent to a validated content export.
Fragmented ordinary streams are tested, whereas general fragmented MFT
traversal is refused. Extension records, attribute lists, sparse/compressed/
encrypted recovery, GPT, complete directory reconstruction, and unknown
payload discovery are not demonstrated.

The team confirmed baseline-file hash matches on November 25.
The 50 new samples share one handcrafted metadata layout.
Recovery-corruption testing was limited to a single-bit substitution.
The separate 100-image detection campaign measured fixed-rule false-positive
and evasion rates, including misses for a reversible marker transformation.
It did not measure unknown-content discovery across independent real volumes.
No benchmark established general extraction throughput. These limits prevent
broad security or compatibility claims even though all included tests passed.

\section{Future Hardening}
The first priority is a broader independent validation corpus: multiple
Windows-formatted images, varied cluster sizes and fragmentation, and
original hashes collected before extraction. A second priority is shared
strict record validation across both CLIs, followed by bounded MFT-run and
attribute-list support. Those changes should be tested against independent
tools and retained failure cases.

Signed manifests would strengthen provenance. Mutation testing and fuzzing
would exercise parser defenses. Separate real-filesystem copy-and-write
tests should measure payload loss. Encryption alone would not resolve these
evidence and parsing limitations.
